\section[Objectifs et contraintes]{Nos objectifs et contraintes pour un
  environnement de dé\-ve\-lop\-pe\-ment et d'observation d'AC}

Nous cherchons à définir et à réaliser un environnement de
développement et d'observation de ces systèmes. %
Une des motivations principales de ce travail réside dans le désir de
voir~(et montrer) les séquences animées générées par un AC\@. %
En effet s'il est aisé de construire des AC en vue de donner des
instantanés photographiques d'états successifs, il est plus difficile de
les visualiser en continu filmique. %
Bien qu'il existe des outils\footnote{%
  \mysoftcite[Cellsim]{cellsim}, %
  \mysoftcite{scamper}, %
  \mysoftcite[CA~LAB]{ca_lab}~\cite[]{Rucker89}, %
  \mysoftcite{xlife}, %
  \mysoftcite[HODGE-C]{hodge_c}, %
  \mysoftcite{cellang}. %
  Les numéros en indice supérieur entre crochets renvoient au répertoire des
  outils, donné en appendice.} %
de spécification d'AC et d'observation de leur dynamique, aucun de
ceux dont nous connaissons l'existence ne répond
pleinement aux contraintes que nous allons définir. %
Nous souhaitons également que cet outil permette immédiatement la
recherche et la création de nouveaux AC, puis la constitution d'un
catalogue d'AC connus. %
Enfin nous espérons qu'il serve à l'identification d'AC pouvant servir
de primitives dans un futur langage basé sur la combinaison d'AC. %
Les qualités recherchées sont la rapidité, l'extensibilité, ainsi que la
facilité de spécification et d'interaction\footnote{%
  La \aconcept{rapidité} peut être définie comme le rendement d'une
  implémentation   pour une architecture donnée. %
  Dans ce contexte, et pour une architecture cible séquentielle, la
  rapidité s'exprime comme le nombre d'instructions
  élémentaires de la machine cible utilisée pour simuler une
  instruction de la machine source. %
  Cette définition est néanmoins trop générale, en effet, si nous
  proposons dans le cadre de ce travail des moteurs cellulaires à
  rendement constant~(\ie indépendant de l'espace initial et de la
  règle de transition), il est possible de concevoir des moteurs
  cellulaires à rendement variable: \mysoftcite{xlife} découpe
  l'espace cellulaire en blocs et permet le calcul sur des
  configurations cellulaires virtuelles de grandes tailles;
  \cite{Gosper84} décrit un algorithme basé sur une mémoire cache
  asynchrone des configurations cellulaires~(le cache contient un
  treillis de morceaux d'espace-temps cellulaire). %

  L'extensibilité peut, dans notre contexte, se ramener à la
  conception d'un ensemble d'objets constituant à la fois la boîte à
  outils et les éléments de base de la construction des expériences
  cellulaires. %
  La recherche de cette qualité peut entrer en conflit avec la
  recherche de la rapidité, par exemple les variations de paramètres
  sur les moteurs cellulaires à rendement constant que nous avons
  réalisés passent par une phase de macro-génération de code source.}. %


\subsection*{Le cadre système, matériel, et logiciel:~les performances}

Notre environnement de développement est Unix et X11, notre
architecture cible est une machine standard de~10~à~100~mips
et~8~à~64 \astrangeterm{mégabytes}{%
  Nous emploierons indifféremment les termes octet et byte, ainsi que
  leurs dérivés: kilobyte, mégaoctet ou gigaoctet.} %
de mémoire utile. %
Pour un AC bidimensionnel composé de \asquare{256}
cellules de~2~bits\footnote{%
  Avec un voisinage de Moore~(9~bits de voisinage), 2~bits d'états
  donnent une table de taille 512~kilobit~($2^{2 \times  9} \times  2$), ce qui
  est très loin des limites que nous nous sommes fixés, mais le passage
  à trois bits de voisinage donne une table de taille
  48~mégaoctets~($\frac{2^{3 \times  9} \times  3}{8}$) qui se trouve à la limite
  des ressources mémoires considérées. %
  Nous insistons cependant sur le fait que rien n'interdit en pratique
  l'usage de telles tables, même si la connaissance de la
  distribution des configurations locales devient nécessaire pour
  prédire les performances du système constitué du triplet
  \emph{table} de transition, \emph{configuration} initiale,
  \emph{performance} intrinsèque du système de gestion mémoire de la
  machine supportant l'application.} %
avec un voisinage de Moore, le système doit être capable de calculer
une dizaine d'états par seconde~(10~Hz)\footnote{%
  D'un point de vue \aquote{esthétique}, cette fréquence correspond à
  un minimum subjectif permettant de restituer l'illusion du mouvement
  continu~(les images télévisuelles ont une fréquence de~25~Hz, les
  images cinématographiques modernes une fréquence de~24~Hz, les
  images cinématographiques originelles une fréquence de~14 à~20~Hz,
  le passage à des fréquences plus élevées est dicté par des
  considérations techniques liées aux besoins du parlant pour le cinéma
  et à la fréquence du secteur alternatif pour la télévision). %
  Cette fréquence est donc élevée, car d'un point de vue
  \aquote{système} elle correspond~(pour notre espace de référence
  $256^2$) à un débit de~640~kilo \bialt{opérations}{cellules} par
  seconde, ce qui pour une machine hypothétique à~20~mips, correspond
  à~32~instructions de la machine cible pour une instruction de la
  machine source~(une mise à jour cellulaire).}. %
Le coût lié à la visualisation est plus difficile à évaluer, mais la
chaîne de traitement complète~(production, extraction, visualisation)
doit permettre l'observation d'un film à une fréquence supérieure
à~1~Hz\footnote{%
  Eisenstein, Gance ou Godard par leur utilisation occasionnelle de
  l'alternance de plans à une telle fréquence illustrent le procédé
  esthétique. %

  Cette fréquence correspond par ailleurs à un impératif ergonomique
  subjectif concernant le \aquote{temps de réponse} d'une application,
  elle borne la taille des espaces accessibles à un mode de
  visualisation séquentielle lorsque la contrainte de l'illusion du
  mouvement est levée. %
  
  La perte de l'illusion du mouvement, n'est pas contradictoire à
  l'étude du mouvement. %
  Le physiologiste Marey étudie le vol des oiseaux:~son %
  \aquote{fusil photographique}~(1882) et son %
  \aquote{Chronophotographe à pellicule}~(1888) %
  servent à l'analyse photographique du mouvement. %
  Bien que considéré comme précurseur du cinématographe,
  il ne se préoccupe pas de restituer l'illusion du mouvement. %
  En cinéma d'animation, la \anewterm{pixillation} utilise sciemment
  la séquence contre l'illusion. %
  \begin{quote}
    ``L'animation n'est pas l'art des dessins qui bougent, mais l'art des
    mouvements dessinés. %
    Ce qui ce passe entre chaque image a beaucoup plus d'importance que
    ce que l'on voit sur chaque image.'' \textsc{N. McLaren}.{\hfill$\cdots/\cdots$}%
  \end{quote}
  En 1941 dans ``L'évolution créatrice'', Bergson
  défend l'idée que ``le mécanisme de notre connaissance usuelle est
  de nature cinématographique''~\cite[page 305]{Bergson89}. %
  Pour lui ``la forme n'est qu'un instantané pris sur une
  transition'', seul est réel le changement continuel de forme. %
  Il place cette métaphore au cœur d'une réflexion historique qui
  remonte à Zénon d'Élée. %

  Dans son acharnement à défendre ``L'antériorité ontologique du
  continu sur le discret'', \cite{Thom92} va jusqu'à proposer
  une ``version informatisée de la preuve cartésienne de
  l'existence de Dieu'' déclarant contradictoire la finitude du nombre
  de neurones cérébraux et la pensée de l'infini. %
  Sans prétendre trancher une question aussi débattue, nous pouvons
  néanmoins suggérer que l'étude des AC propose une approche à rebours
  de cette question, la recherche pratique de l'immanence
  phénoménologique du continu remplaçant les tentatives de
  démonstration de son antériorité ontologique. %

  Sur notre configuration de développement~(un~486DX2 avec~32~mégaoctets
  de mémoire, le système d'exploitation Linux et X11 pour la
  visualisation) la fréquence typique pour une chaîne comprenant la
  génération sur un voisinage de Moore~(toujours pour un espace de
  référence $256^2$) et l'extraction d'un plan mémoire~(constitué d'un
  des deux bits d'état) dans un premier processus, pipé sur un processus
  dumpant les plans cellulaires vers le serveur X11~(tournant sur la
  même machine) est supérieure à~10~Hz. %
  Le découplage des différentes opérations de la chaîne permet
  d'obtenir une fréquence de génération supérieure à~15~Hz~(module de
  génération redirigé sur \texttt{/dev/null}) et une fréquence de
  visualisation monochrome~(pour un seul plan de bits cellulaire)
  supérieure à~30~Hz, (mesuré via un fichier historique
  de~8~mégaoctets pour~1024~générations). %
  Ces mesures s'entendent en temps écoulé réel, pour un système
  autrement au repos, les mesures de performance de visualisation ne
  constituent évidemment pas une référence à notre travail, mais elles
  permettent de considérer, au regard des mesures de performance du
  moteur cellulaire, que celui-ci atteint le débit plafond d'un moteur
  à rendement constant, \ie la meilleure performance possible est
  toujours garantie dans le pire des
  cas~(espace cellulaire et portion d'espace-temps saturé).}. %
Nous avons ainsi réalisé un noyau d'objets permettant la construction
d'expériences, sous la forme d'une librairie C, interfacée à
\asoft{Tcl}\footnote{%
  Tool Command Language.}%
~\cite[]{Ousterhout90,Ousterhout91}. %

\subsection*{Spécification:~la description des expériences cellulaires}

La spécification concerne les réseaux d'interconnexions, les fonc\-tions
de transitions, et le séquencement des opérations applicables sur
l'espace. %
Les fonc\-tions de transitions sont écrites en~\asoft{C}, et sont dotées
d'une convention d'interface prévue pour permettre une indépendance
des paramètres à leur localisation dans la mémoire d'états d'une
cellule et la réutilisation dans des fonctions de plus haut niveau. %
Les réseaux d'interconnexions sont définis et créés en~\asoft{Tcl} à
partir des primitives de la librairie. %
Les espaces sont créés en~\asoft{Tcl}. %
Le séquencement est défini par un programme~\asoft{Tcl}. %

\subsection*{Interaction:~l'interface humaine}

L'interaction se fait soit par l'intermédiaire d'un shell~\asoft{Tcl}, soit
par une interface X11 basé sur~\asoft{Tk}. %

\subsection*{Visualisation:~les outils graphiques}

La visualisation se fait par l'intermédiaire de fenêtres X11
présentant des séquences d'espaces monochromes~(plans de bits de
l'espace cellulaire) ou couleur~(tranches de bits de l'espace
cellulaire). %
Les variables macroscopiques~(éléments statistiques) sont visualisées
par l'intermédiaire d'objets X11 spécialisés~(plotters). %

\mydumpfull[htb]{casNew}%
{Une interface humaine expérimentale pour le Workbench}%
{Une interface humaine expérimentale et son dumper}%
{Une interface X11 réalisée avec
  \protect\mysoftcite{Tk}~\cite[]{Ousterhout90,Ousterhout91}. %
  Le dumper présente une image typique de la transformation d'un espace
  initial aléatoire de densité\,\Sundemi par la
  règle~\arule{Anneal}~\cite[]{Vichniac86,Toffoli87}\footnote{%
    \mygetrule[phrase]{anneal}} %
  Le dumper présente également différents compteurs liés aux temps
  propre de la règle, du séquenceur, et au pool d'espaces de l'animation. %
}

\section{Les outils existants:~pas de standard}

Une recherche aussi exhaustive que possible des outils publics
existants nous a convaincu de la nécessité d'une utilisation critique
de certains de ces outils, en particulier dans le domaine de la
visualisation, associé à une réalisation expérimentale ad hoc
d'éléments d'architecture de calcul cellulaire. %

Les outils de développement et de visualisation d'automates
cellulaires sont peu nombreux\footnote{%
  Bien que les réalisations de \arule{Life} soit
  innombrables, et les packages cellulaires plus ou moins généraux très
  nombreux, leurs réalisations minimales posant peu de problèmes, nous ne
  prenons en compte que les packages dont il existe des annonces
  explicites.}. %
Le plus connu d'entre eux est sans doute \mysoftcite{cellsim};
\mysoftcite{scamper} et \mysoftcite[CA~LAB]{ca_lab} sont également des
outils disponibles depuis 1989. %
Il existe aussi des outils spécialisés pour un modèle d'automate
cellulaire. %
\mysoftcite{xlife} fournit une réalisation d'exécution très rapide du
jeu de la vie~\cite[]{Berlekamp82}, \mysoftcite[HODGE-C]{hodge_c} une
réalisation de la machine hodge-podge de
Gerhard~\&~Schuster~\cite[]{Dewdney88}, une classe d'AC inspirée par
les réactions chimiques autocatalytiques, telle que la réaction de
Beluso\mbox{v-Z}habotinski. %
Enfin, \mysoftcite{cellang}, d'inspiration plus architecturale,
propose un langage de description d'AC~\cite[]{Eckart92a}. %
\nocite{Eckart92b}

\section[Le calcul et la visualisation]%
{Un nouvel environnement:~exploration par le calcul et la visualisation}

Nous avons donc réalisé un nouvel environnement\footnote{%
  Ci-après désigné par le terme \anewterm{workbench}.} %
de développement et de visualisation d'automates cellulaires. %

Nous destinons ce workbench tant à l'exploration de la dynamique
des automates cellulaires~(AC), qu'à la construction et à l'évaluation
d'architectures et d'algorithmes SIMD\@. %
De plus nous proposons de considérer la réalisation d'algorithmes SIMD
sur machines SISD non pas comme un palliatif à l'absence de machines
SIMD, mais comme un mode d'utilisation particulier des ressources
disponibles sur les machines SISD modernes~(50~Mips, 32~mégaoctets). %

La programmation des calculateurs massivement parallèles réels et
virtuels pose le double problème du pouvoir d'expression et de
l'efficacité d'un langage au regard d'une architecture. %
Ce problème se pose également pour les architectures SISD, mais l'on
peut considérer qu'il concerne désormais l'adéquation des langages aux
applications plutôt qu'aux architectures. %
Optant pour une approche résolument bottom-up de la recherche des
primitives, nous proposons d'utiliser les AC comme composants
élémentaires d'une combinatoire visant à la construction d'objets
\bialt{cellulaires}{SIMD} arbitrairement complexes. %
Ces objets se divisent en deux classes épistémologiques. %
Nous les nommerons objets \aconcept{analytiques} et objets
\aconcept{synthétiques}. %
Les objets analytiques sont typiquement issus de l'application
itérative d'une fonction simple sur un espace cellulaire initial
aléatoire. %
Ils recoupent intuitivement le champ d'étude traditionnel des AC\@. %
Les objets synthétiques sont issus de la sélection d'une fonction
localement universelle, et de son application contrôlée~(globalement
paramétrée) sur un espace initial soigneusement configuré. %
Ils correspondent aux architectures des machines SIMD\@. %
De l'étude des objets analytiques nous espérons tirer des
enseignements permettant à terme de considérer les processus
morphogénétiques comme des éléments de programmation. %
De la construction d'un outillage SIMD nous attendons les moyens de
contrôle nécessaires à l'étude des objets analytiques. %
Cet objectif, certes ambitieux, repose précisément sur la réalisation
par \aconcept{lookup-tables} des fonctions de bas niveau manipulables
et combinables au sein de l'environnement. %

Notre étude est organisée autour de la présentation des structures de
données choisies pour le workbench; ces structures de données
représentent \aconcept{l'espace}, \aconcept{le temps}, et
\aconcept{les lois}:~objets élémentaires du workbench. %
Nous proposons des exemples d'opérations applicables sur ces objets. %
Nous proposons également de considérer la réalisation de fonctions par
lookup-tables comme une structure de données fondamentale, tant pour
des raisons d'efficacité, que de généralité. %

Nous considérons la visualisation comme un service primordial, qui
propose l'usage du workbench en tant que laboratoire de manipulation
d'espace cellulaire~(l'espace, mais aussi les fonctions et les
dynamiques). %
Nous présentons plusieurs exemples de visualisations principalement
associées à des objets analytiques. %
Nous montrons comment peut s'effectuer le passage d'une méthode
naturaliste à une méthode artefactuelle en proposant une réalisation
par lookup-table d'une architecture SIMD spécifique, et une méthode de
transformation d'une lookup-table quelconque en séquence de contrôle
pour cette machine. %

Constatant l'absence d'outils adaptés aux contraintes de la
visualisation des espace-temps cellulaires\footnote{%
  La visualisation de volume de données dans un format inadapté
  conduit à la saturation inutile des capacités mémoire de la machine
  d'accueil d'un calcul cellulaire.}, %
constatant également l'existence d'espace-temps cellulaires binaires a
priori toriques, nous avons conçu, réalisé et testé une bibliothèque
de manipulation d'espace\footnote{%
  Coupe, séquences de coupes sur des axes de projection variables, %
  réduction intégrale sur un axe arbitraire pondéré par la distance à
  la surface de projection pour obtenir des effets de transparence.} %
cellulaires aux fins de visualisation. %
Sur cette base, nous avons par exemple pu tester l'idée d'une
visualisation relativiste des espace-temps:~%
nous avons projeté le cône de la vitesse de la lumière\footnote{%
  La vitesse de la lumière d'un espace-temps cellulaire est le rayon
  du voisinage de la famille de règle à laquelle appartient la règle
  génératrice de l'espace-temps.} %
d'un espace-temps~2D sur un espace~2D, résumant ainsi sur un seul
espace~2D la dynamique de la règle visualisée.

\mydumpfull[p]{superglider}%
{Un glider superluminique}%
{Un glider subluminique et un glider superluminique}%
{Nous présentons dans cette figure la découverte
  d'un \aquote{glider superluminique}. %
  Cette série d'espaces de configurations cellulaires est une séquence
  de tranches d'espace-temps non alignées sur l'axe du temps. %
  Malgré la distorsion introduite par ce décalage, la règle
  génératrice de cette portion d'espace-temps est
  parfaitement identifiable:~\arule{Life}\footnote{%
    Les flèches en haut à droite des deux premières photos pointent
    une configuration cyclique de~\arule{life} reconnaissables.}. %
  A première vue, la conséquence de l'utilisation d'un axe non
  standard pour la projection de la dynamique de~\arule{Life} est
  d'ordre purement spatial\footnote{%
    Cette appréciation subjective n'est applicable qu'à l'expérience
    de la vision du \aquote{film}.}, %
  exception faite du décalage introduit par le désaxage temporel. %
  Le carré sur les photos pointe vers un objet \arule{Life} de type
  glider, doté d'une vitesse de déplacement subluminique. %
  L'observation de l'espace-temps contenant la portion présentée au
  moyen de films à haute fréquence\footnote{%
    Supérieure à la fréquence télévisuelle.} %
  nous a permis d'identifier à la vue un objet paradoxal, pointé par
  un cercle sur les photos:~un glider superluminique.}

\noindent La figure~\mygetdump{superglider} présente un des usages
possible de notre librairie de manipulation d'espaces dans le
contexte du workbench:~%
nous y présentons la découverte d'un \aconcept{glider
  superluminique}\footnote{%
  nous emploierons le terme glider~(planeur) de manière générique, \ie
  pour désigner des objets cellulaires périodiques, à un décalage spatial près.}, %
en déplacement à une vitesse apparente double de la
vitesse de la lumière. %
Nous attendions évidemment de cette expérience des aberrations
temporelles, des illusions ou mirages inex\-plicables dans le contexte
d'une projection parallèle à l'axe du temps. %
Néanmoins les premiers résultats faisaient plutôt apparaître une
relative invariance de la perception visuelle subjective associée à la
règle. %
La procédure utilisée produit un déplacement apparent des
configurations, mais les objets se comportent par ailleurs, c'est à
dire relativement à l'espace apparent pour l'expérimentateur, de
manière subluminique ou luminique usuelle. %
Ce glider, le seul que nous ayons trouvé~(dans l'espace-temps
englobant la portion présentée) était donc prévisible, puisqu'il
correspond à un glider en mouvement corrélé au front de coupe, mais
néanmoins totalement surprenant. %
Beaucoup moins prévisibles nous apparaissent les conséquences
générales d'observations exploitant systématiquement des projections
d'espace-temps non parallèles à l'axe du temps.

\myfigfullR[htb]{.75}{timepyrcub}%
{Tranches d'espace-temps et pyramides relativistes}%
{Schémas de réalisation de tranches d'espace-temps et de pyramides relativistes}%
{Le cube, sur ces schémas, représente une portion d'espace-temps pour
  un AC~2D. %
  L'axe du temps est indiqué par la flèche verticale, %
  les pointillés larges sont des arêtes de plans de coupe. %
  Le schéma de gauche illustre la procédure utilisée pour réaliser les
  séquences de coupes présentées en figure~%
  \mygetdump{superglider}, %
  \mygetdump{tsum-128x3-diag}, %
  et \mygetdump{anneal-256x3-diag}. %
  Les coupes sont calculées en \aquote{distance de Manhattan}, %
  Nous avons conçu un calcul d'intersection basé sur une
  version~3D de l'algorithme de tracé de ligne de Bresenham\footnote{%
    Nous avons modifié pour cet usage la version de Bob
    Pendleton~(1992) dérivée de~\cite{Heckbert90}.}. %
  Le schéma de droite illustre la procédure utilisée pour réaliser les
  projections relativistes\footnote{%
    La projection est construite par \aconcept{cercles de voisinage}
    remontant le temps, dont le rayon détermine l'âge, éventuellement
    associés à une translation spatiale figurée par la flèche horizontale
    pleine.} %
  présentées en figures~%
  \mygetdump{life-256x3-rel}, %
  \mygetdump{life-256x3-relmov}, %
  \mygetdump{tsum-128x3-relmovedge}, %
  }

\mydumpfull[p]{tsum-128x3-diag}%
{Une séquence de tranches d'espace-temps}%
{Une séquence de tranches d'espace-temps de la règle~\arule{Qmean}}%
{La règle~\arule{Qmean} est présentée, ainsi que les autres règles
  utilisées pour la synthèse des images réalisées
  en~\mygetref[name]{the-rules}. %
  La figure~\mygetfig{timepyrcub} illustre la procédure de réalisation
  de cette séquence. %
  D'autres visualisations associées à la règle~\arule{Qmean} sont
  présentées par les figures~%
  \mygetdump{tsum-grey}, %
  \mygetdump{tsum}, %
  \mygetdump{tsum-ray-one} %
  et \mygetdump{tsum-ray-all}.}

\mydumpfull[p]{life-256x3-rel}%
{Projection relativiste fixe}%
{Une projection relativiste pour une cellule fixe}%
{Une projection relativiste~\mygetfig[paren]{timepyrcub} d'un
  espace-temps de \arule{life} pour une cellule de référence fixe,
  réalisée sur un espace-temps~\acube{256}, pour une densité
  initiale\,\Sundemi. %
  La cellule de référence est au centre de la projection, %
  celle-ci apparaît donc au sein d'un espace \aquote{âgé},
  composé en plus grand nombre d'objets point-fixe et %
  de densité moyenne  plus faible %
  que l'espace situé à proximité des bords de la projection, %
  de densité moyenne plus élevée et peuplé d'objets plus jeunes\footnote{%
    Que l'on peut cependant qualifier d'archaïques du point de vue de
    l'origine de la projection.}. %
  Il est possible d'accentuer l'âge des objets lointains en diminuant la
  vitesse de la lumière utilisée lors de la projection. %
  Un tel découplage du rayon de voisinage accroît le problème
  d'interprétation des projections produites.}

\mydumptwothirds[htb]{life-256x3-relmov}%
{Projection relativiste mobile}%
{Une projection relativiste à cellule mo\-bile}%
{Cette projection diffère de la figure~\mygetdump{life-256x3-rel} par
  l'ajout d'une translation appliquée à la cellule de référence.}

\mydumpfull[t]{anneal-256x3-rel}%
{}%
{Projection relativiste sur les 3~axes d'un cube d'espace-temps}%
{Cette figure regroupe trois projections relativistes de la
  règle~\arule{Anneal} réalisées sur une même trajectoire, %
  pour l'axe temporel et les deux axes nord-sud et est-ouest. %
  L'origine du temps est visible dans le haut de la projection pour
  ces deux axes.}

\myfigfullR[htb]{.75}{transcub}%
{Transparence globale et par plan de bits}%
{Schéma de réalisation de réductions par transparence}%
{Comme pour la figure~\mygetfig{timepyrcub}, les cubes représentent une
  portion d'espace-temps pour un AC 2D. %
  La flèche pointillée illustre une traversée orthogonale non parallèle
  à l'axe du temps, lors d'une telle traversée on peut soit réaliser un
  film\footnote{%
    A défaut de films, incompatible avec le support papier, nous
    présentons des séquences associées à une traversée non
    orthogonale~\mygetfig[paren]{timepyrcub}.}, %
  soit une réduction intégrale~(une somme de la valeur de toutes
  les cellules rencontrées). %
  Pour les réductions intégrales destinées à la production d'images,
  la valeur de la cellule est pondérée par la distance parcourue. %
  La figure~\mygetdump{tsum-128x3-trans-byte} présente une telle
  image. %
  Pour les transparences non orthogonales, les effets de bord de la
  méthode utilisée~(distance de Manhattan) posent un problème
  d'interprétation des déformations. %
  Les quatre cubes à droite symbolisent l'opération de séparations d'un
  espace-temps de cellules à quatre bits en quatre espace-temps de
  cellules à un bit. %
  La figure~\mygetdump{tsum-128x3-trans-bits} présente des
  transparences sur trois axes pour cinq \aquote{plans} de bits
  de~\arule{Qmean}, %
  La figure~\mygetdump{life-256x3-trans-bits} pour l'unique plan d'un
  espace-temps de~\arule{life}.}

\mydumpfull[htb]{anneal-256x3-diag}%
{}%
{Une séquence de tranches d'espace-temps de la règle~\arule{Anneal}}%
{Les premières coupes correspondent à un temps proche du temps initial,
  mais chaque tranche présente des caractéristiques archaïques plus ou
  moins marquées.}

\mydumpfull[htb]{tsum-128x3-faces}%
{}%
{Faces du cube d'espace-temps de~\arule{Qmean}}%
{Cette figure présente les six faces d'un cube d'espace-temps de la
  règle~\arule{Qmean}. %
  Les deux faces situées dans l'axe du temps, ainsi que leur âges
  sont clairement identifiables.}

\mydumpfull[p]{tsum-128x3-trans-bits}%
{}%
{Une transparence éclatée par plan de bits}%
{Les cinq bits de poids fort de la règle~\arule{Qmean} sont utilisés
  respectivement pour réaliser une
  transparence~\mygetfig[paren]{transcub} sur les trois axes de
  l'espace-temps. %
  Ces séquences sont à comparer à la
  figure~\mygetdump{tsum-128x3-trans-byte} utilisant la totalité des
  bits de chaque cellule pour la transparence.}

\mydumpfull[p]{tsum-128x3-trans-byte}%
{}%
{Une transparence utilisant tout l'espace de valeur d'une cellule}%
{Cette transparence fournit une image plus habituelle d'une
  dynamique cellulaire considérée comme un voxmap\footnote{%
    Le voxel est la version 3D du pixel.} %
  Des expériences antérieures menées avec des packages de visualisation
  scientifique \aquote{classiques}\footnote{%
    Au sens ou ils travaillent nécessairement en flottants, ou considèrent
    des données dont la localisation n'est pas implicite
    et en tout cas ne sont pas adaptés à un travail au niveau du bit. %
    Nous avons recherché, porté et expérimenté les outils suivants:~%
    \mysoftcite[XDataSlice]{XDS} et son format %
    \mysoftcite{HDF}, %
    \mysoftcite{VIS-5D}~\cite[]{Hibbard94}, %
    \mysoftcite{viewit} et %
    \mysoftcite{LIC}~\cite[]{Cabral93}.} %
  donnent pour ce genre de visualisation des résultats proches en terme
  d'images produites\footnote{%
    Mais pas en termes de ressources consommées pour la production. %
    Un autre inconvénient de ces packages, étant les erreurs
    introduites par la conversion flottante intermédiaire lorsque
    l'espace 2D de transparence produit est destiné à la réalisation
    de mesures statistiques, autant que d'images.} %
  de ceux présentés ici, %
  mais l'usage de nos fonc\-tions de manipulation d'espaces cellulaires
  est nécessaire\footnote{%
    D'un point de vue pratique, la conversion d'un bitmap
    tridimensionnel en voxmap de flottants, éventuellement reconverti
    en voxmap de bytes, est en effet parfaitement inadéquate.} %
  à la pro\-duc\-tion des transparences par plan de bits de
  la figure~\mygetdump{tsum-128x3-trans-bits}. %
  La comparaison des images produites par les méthodes plan de bits
  \vs voxel, suggère l'étude des relations entre les deux
  formulations, et une possible généralisation de la première à
  l'analyse de voxmaps.} %

\mydumpfull[htb]{anneal-256x3-trans-bits}%
{}%
{Une transparence par plan de bits pour~\arule{Anneal}}%
{La méthode de production commentée en
  figure~\mygetfig{transcub} est utilisée 
  pour la règle~\arule{Anneal}. %
  Le nombre d'états de~\arule{Anneal} étant égal à deux, la méthode
  éclatée se confond avec la méthode standard.}

\mydumpfull[htb]{life-256x3-trans-bits}%
{}%
{Une transparence par plan de bits pour~\arule{Life}}%
{La méthode de production commentée en
  figure~\mygetfig{transcub} est utilisée 
  pour la règle~\arule{Life}}

\mydumpfull[htb]{tsum-128x3-relmovedge}%
{}%
{Une détection de contour appliquée sur une projection relativiste}%
{L'accumulation de transformations renforce le problème de
  l'interprétation des images produites. %
  Ces figures illustrent les limites rencontrées en combinant une
  détection de contour une à projection relativiste à cellule de
  référence mobile, appliquée sur les trois axes de projection d'un
  cube d'espace-temps de la règle~{Qmean}.}

L'état actuel du développement du workbench n'a pas encore permis
l'expérimentation de constructions complexes d'objets mixtes. %
La mesure de variables d'états macroscopiques, par exemple, utilise
actuellement des outils externes, au lieu d'algorithmes SIMD adaptés à
la nature discrète et finie des espaces cellulaires. %

%%
%% Local Variables:
%% mode: auto-fill
%% iso-accents-minor-mode: nil
%% TeX-command-default: "mytex"
%% TeX-master: "top"
%% Local IspellParsing: "latex-mode ~latin1"
%% Local IspellPersDict: "~/.ispell_lisp"
%% LocalWords: "d\'evelop pement d'AC cellsim scamper CA LAB ca lab xlife hodge"
%% LocalWords: "HODGE-C cellang impl\'ementation d'espace-temps mips m\'egabytes"
%% LocalWords: "byte kilobyte m\'egaoctet gigaoctet bidimensionnel Moore kilobit"
%% LocalWords: "Eisenstein Gance Godard Marey Chronophotographe pixillation ur"
%% LocalWords: "McLaren Bergson Z\'enon d'\'El\'ee finitude ph\'enom\'enologique DX Linux"
%% LocalWords: "pip\'e dumpant dev null monochrome interfac\'ee Tcl Tool Command htb"
%% LocalWords: "Language s\'equencement casNew Workbench dumper Tk Anneal shell ad"
%% LocalWords: "s\'equenceur monochromes plotters hoc packages hodge-podge Gerhard"
%% LocalWords: "Schuster Beluso v-Z macro-g\'en\'eration habotinski workbench SIMD"
%% LocalWords: "SISD bottom-up lookup-tables manipulables artefactuelle NEW tsum"
%% LocalWords: "superluminique subluminique luminique timepyrcub relativiste rel"
%% LocalWords: "diag anneal Manhattan Bresenham Pendleton relmov relmovedge XDS"
%% LocalWords: "Qmean the-rules tsum-grey tsum-ray-one tsum-ray-all paren voxmap"
%% LocalWords: "est-ouest transcub trans-byte trans-bits voxel XDataSlice HDF"
%% LocalWords: "viewit LIC bitmap bytes versus voxmaps nord-sud name"
%% LocalWords: "lookup-table relativistes glider point-fixe"
%% LocalWords: "espace-temps l'espace-temps superglider d\'esaxage"
%% End:
